% GENERATED FILE. Edit canonical lessons, then run npm run generate:books.
% canonical-source-hash: fnv1a64:20010f4f04f91bfe
% canonical-lessons: GE-C85-eis, GE-C85-huegel, GE-C85-wueste, GE-C85-pferd, GE-C85-kuh

\chapter{Ice, Hill, Desert, Horse, Cow}
\label{ch:ge-ice-hill-desert-horse-cow}
\hlchaptermodality{85}

\begin{chapteropening}
\textbf{By the end of this chapter:} \emph{I can say ice, a hill, a desert, a horse and a cow.}

Complete the last lesson of chapter 85: I can say ice, a hill, a desert, a horse and a cow.
\end{chapteropening}

\section[das Eis]{das Eis — ice; ice cream}

\label{lesson:GE-C85-eis}



\begin{quote}
Before the new one: say the German for a pond, then the German for fire.
\end{quote}

\subsection*{You'll want to know: das Eis}
\textbf{das Eis} — \textquotedblleft{}ice; ice cream\textquotedblright{}.

\begin{cousinweb}
\textbf{Eis} $\leftarrow$ Old High German \emph{īs} $\leftarrow$ Proto-Germanic \emph{*īsą} — inherited, the same word as English \emph{ice}. (EtymDB 2.1) \textquotedblleft{}Ice cream\textquotedblright{} is a later sense.
\end{cousinweb}

\begin{grammarlens}[title={What you've built}]
One more word for this chapter.
\end{grammarlens}

\subsection*{Your turn}
\begin{itemize}
\raggedright
  \item \emph{Say it:} \emph{das Eis}
  \item \emph{Say it:} \emph{das Eis}, once more
  \item \emph{Say it:} say \emph{das Eis}
  \item \emph{Recall:} say the German for a pond, then the German for fire, then say \emph{das Eis} again
\end{itemize}

\begin{tcolorbox}[breakable,colback=teal!4,colframe=teal!35!black,title={Before you move on}]
What does das Eis mean? (\textquotedblleft{}Ice; ice cream\textquotedblright{}.) Say it once more.
\end{tcolorbox}
\section[der Hügel]{der Hügel — a hill}

\label{lesson:GE-C85-huegel}



\begin{quote}
Before the new one: say the German for fire, then the German for ice; ice cream.
\end{quote}

\subsection*{You'll want to know: der Hügel}
\textbf{der Hügel} — \textquotedblleft{}a hill\textquotedblright{}.

\begin{cousinweb}
\textbf{Hügel} spread into standard German from central German dialects. It is a diminutive of Middle High German \emph{houc}, \textquotedblleft{}hill\textquotedblright{}, related to the word for \textquotedblleft{}high\textquotedblright{}; compare Old Norse \emph{haugr}, \textquotedblleft{}mound\textquotedblright{}.
\end{cousinweb}

\begin{grammarlens}[title={What you've built}]
One more word for this chapter.
\end{grammarlens}

\subsection*{Your turn}
\begin{itemize}
\raggedright
  \item \emph{Say it:} \emph{der Hügel}
  \item \emph{Say it:} \emph{der Hügel}, once more
  \item \emph{Say it:} say \emph{der Hügel}
  \item \emph{Recall:} say the German for fire, then the German for ice; ice cream, then say \emph{der Hügel} again
\end{itemize}

\begin{tcolorbox}[breakable,colback=teal!4,colframe=teal!35!black,title={Before you move on}]
What does der Hügel mean? (\textquotedblleft{}A hill\textquotedblright{}.) Say it once more.
\end{tcolorbox}
\section[die Wüste]{die Wüste — a desert}

\label{lesson:GE-C85-wueste}



\begin{quote}
Before the new one: say the German for ice; ice cream, then the German for a hill.
\end{quote}

\subsection*{You'll want to know: die Wüste}
\textbf{die Wüste} — \textquotedblleft{}a desert\textquotedblright{}.

\begin{cousinweb}
\textbf{Wüste} $\leftarrow$ Old High German \emph{wuostī} $\leftarrow$ Proto-Germanic \emph{*wōstinjō}, \textquotedblleft{}wasteland\textquotedblright{} — inherited. (EtymDB 2.1) English \emph{waste}, which came through French from Latin \emph{vastus}, is a cousin.
\end{cousinweb}

\begin{grammarlens}[title={What you've built}]
One more word for this chapter.
\end{grammarlens}

\subsection*{Your turn}
\begin{itemize}
\raggedright
  \item \emph{Say it:} \emph{die Wüste}
  \item \emph{Say it:} \emph{die Wüste}, once more
  \item \emph{Say it:} say \emph{die Wüste}
  \item \emph{Recall:} say the German for ice; ice cream, then the German for a hill, then say \emph{die Wüste} again
\end{itemize}

\begin{tcolorbox}[breakable,colback=teal!4,colframe=teal!35!black,title={Before you move on}]
What does die Wüste mean? (\textquotedblleft{}A desert\textquotedblright{}.) Say it once more.
\end{tcolorbox}
\section[das Pferd]{das Pferd — a horse}

\label{lesson:GE-C85-pferd}



\begin{quote}
Before the new one: say the German for a hill, then the German for a desert.
\end{quote}

\subsection*{You'll want to know: das Pferd}
\textbf{das Pferd} — \textquotedblleft{}a horse\textquotedblright{}.

\begin{cousinweb}
\textbf{Pferd} $\leftarrow$ Old High German \emph{pfarifrit} $\leftarrow$ Late Latin \emph{paraverēdus}, \textquotedblleft{}spare post horse\textquotedblright{} (EtymDB 2.1) — borrowed and worn down. English \emph{palfrey} comes from the same Latin word through French.
\end{cousinweb}

\begin{grammarlens}[title={What you've built}]
One more word for this chapter.
\end{grammarlens}

\subsection*{Your turn}
\begin{itemize}
\raggedright
  \item \emph{Say it:} \emph{das Pferd}
  \item \emph{Say it:} \emph{das Pferd}, once more
  \item \emph{Say it:} say \emph{das Pferd}
  \item \emph{Recall:} say the German for a hill, then the German for a desert, then say \emph{das Pferd} again
\end{itemize}

\begin{tcolorbox}[breakable,colback=teal!4,colframe=teal!35!black,title={Before you move on}]
What does das Pferd mean? (\textquotedblleft{}A horse\textquotedblright{}.) Say it once more.
\end{tcolorbox}
\section[die Kuh]{die Kuh — a cow}

\label{lesson:GE-C85-kuh}



\begin{quote}
Before the new one: say the German for a desert, then the German for a horse.
\end{quote}

\subsection*{You'll want to know: die Kuh}
\textbf{die Kuh} — \textquotedblleft{}a cow\textquotedblright{}.

\begin{cousinweb}
\textbf{Kuh} $\leftarrow$ Old High German \emph{kuo} $\leftarrow$ Proto-Germanic \emph{*kūz} — inherited, the same word as English \emph{cow}. (EtymDB 2.1)
\end{cousinweb}

\begin{grammarlens}[title={What you've built}]
One more word for this chapter.
\end{grammarlens}

\subsection*{Your turn}
\begin{itemize}
\raggedright
  \item \emph{Say it:} \emph{die Kuh}
  \item \emph{Say it:} \emph{die Kuh}, once more
  \item \emph{Say it:} say \emph{die Kuh}
  \item \emph{Recall:} say the German for a desert, then the German for a horse, then say \emph{die Kuh} again
\end{itemize}

\begin{tcolorbox}[breakable,colback=teal!4,colframe=teal!35!black,title={Before you move on}]
What does die Kuh mean? (\textquotedblleft{}A cow\textquotedblright{}.) Say it once more.
\end{tcolorbox}
